%% kapa: exam=1 round=21 year=2010 period=1 form=A subject=civil_law subject_ko=민법 seq=8 number=8
%% source: Q-Net 감정평가사 기출문제 (cst00309) artlSeq=5024181 file=extracted/1008/5024181/2010년 제21회 감정평가사 1차시험 기출문제/1교시.hwp
%% conversion: kordoc 4.24.51 -> tools/gen1.py (auto); review=proofread-sonnet; flags=none
\begin{question}{8}
권리의 객체에 관한 설명으로 옳지 않은 것은? (다툼이 있으면 판례에 의함)
\begin{choices}
  \item 물건의 사용이익에 관하여는 과실에 관한 규정을 유추적용할 수 있다.
  \item 건물 그 자체의 효용과 직접 관계가 없는 물건은 건물의 종물이 아니다.
  \item 증감 변동하는 집합동산도 다른 물건과 구별될 수 있도록 특정성을 갖추면, 그 전부를 목적으로 하는 양도담보권이 유효하게 성립할 수 있다.
  \item 종물은 독립한 물건이어야 하지만, 반드시 동산이어야 하는 것은 아니다.
  \item 건물의 일부에는 전세권을 설정할 수 없다.
\end{choices}
\end{question}
\answer{5}
